\section*{Bab \ref{ch:g10:reffunc} --- Fungsi Acuan}

\begin{solution}{exo:g10:reffunc:1}
$f$: \omterm{def:g10:reffunc:affine}{gradien} $3$, titik potong $-2$, naik. \quad
$g$: \omterm{def:g10:reffunc:affine}{gradien} $-\frac12$, titik potong $4$, turun. \quad
$h$: \omterm{def:g10:reffunc:affine}{gradien} $0$, titik potong $7$, tetap. \quad
$k$: \omterm{def:g10:reffunc:affine}{gradien} $-1$, titik potong $0$, turun (dan linear).
\end{solution}

\begin{solution}{exo:g10:reffunc:2}
Melalui $(2,1)$ dan $(5,10)$: \omterm{def:g10:reffunc:affine}{gradien} $m = \frac{10 - 1}{5 - 2} = 3$; lalu
$1 = 3 \times 2 + p$ memberi $p = -5$: $f(x) = 3x - 5$.

Melalui $(-1,4)$ dan $(3,-4)$: \omterm{def:g10:reffunc:affine}{gradien} $m = \frac{-4 - 4}{3 - (-1)} =
\frac{-8}{4} = -2$; lalu $4 = -2 \times (-1) + p$ memberi $p = 2$:
$f(x) = -2x + 2$.
\end{solution}

\begin{solution}{exo:g10:reffunc:3}
$(3.1)^2 < (3.2)^2$: pangkat dua naik pada bilangan positif.

$(-2.7)^2 < (-2.8)^2$: pangkat dua turun pada bilangan negatif dan
$-2.8 < -2.7$.

$\frac{1}{5.1} > \frac{1}{5.2}$: kebalikan turun pada bilangan positif.

$\sqrt{17} > \sqrt{15}$: akar kuadrat naik.
\end{solution}

\begin{solution}{exo:g10:reffunc:4}
$x^2 = 16$: dua penyelesaian, $x = 4$ dan $x = -4$ (garis mendatar
$y = 16$ memotong parabolanya dua kali).

$x^2 \leq 16$: parabolanya berada di bawah garis itu di antara kedua
titik potongnya: $x \in \intcc{-4}{4}$.

$x^2 > 9$: parabolanya tegas berada di atas $y = 9$ di luar
$\intcc{-3}{3}$: $x \in \intoo{-\infty}{-3} \cup \intoo{3}{+\infty}$.
\end{solution}

\begin{solution}{exo:g10:reffunc:5}
Untuk $0 < x < 1$: mengalikan $x < 1$ berulang kali dengan $x$ memberi
$x^3 < x^2 < x$, dan \cref{prop:g10:reffunc:compare} memberi $x < \sqrt x$,
sehingga
\[
x^3 < x^2 < x < \sqrt x .
\]
Periksa dengan $x = 0.25$: $x^3 = 0.015625$, $x^2 = 0.0625$, $x = 0.25$,
$\sqrt x = 0.5$.
\end{solution}

\begin{solution}{exo:g10:reffunc:6}
$\frac1x = 2$ mempunyai penyelesaian tunggal $x = \frac12$ (hiperbolanya
memotong garis mendatar $y = 2$ sekali, pada cabang positifnya).

Untuk $x > 0$: $\frac1x < 2$. Mengalikan dengan $x > 0$: $1 < 2x$, jadi
$x > \frac12$: himpunan penyelesaiannya $\intoo{\frac12}{+\infty}$.

Bila $x < 0$ juga diizinkan: setiap $x$ negatif memenuhi $\frac1x < 0 < 2$,
sehingga himpunan penyelesaian lengkapnya adalah
$\intoo{-\infty}{0} \cup \intoo{\frac12}{+\infty}$.
\end{solution}

\begin{solution}{exo:g10:reffunc:7}
(a) Pada $\intcc{2}{5}$, semua masukannya positif dan pengkuadratan naik:
$4 \leq x^2 \leq 25$.

(b) Pada $\intcc{-3}{-1}$, pengkuadratan turun:
kuadrat terbesarnya berasal dari $-3$: $1 \leq x^2 \leq 9$.

(c) Pada $\intcc{-2}{3}$, \omterm{def:g10:numbers:interval}{intervalnya} memuat $0$, tempat kuadratnya
mencapai \omterm{def:g10:functions:extrema}{minimum} $0$; \omterm{def:g10:functions:extrema}{maksimumnya} adalah yang lebih besar di antara
$(-2)^2 = 4$ dan $3^2 = 9$. Jadi $0 \leq x^2 \leq 9$.
\end{solution}

\begin{solution}{exo:g10:reffunc:8}
Perusahaan pertama: $f(x) = 1.5x + 4$; yang kedua: $g(x) = 2x$. Yang pertama
lebih murah apabila
\[
1.5x + 4 < 2x
\ \Longleftrightarrow\
4 < 0.5x
\ \Longleftrightarrow\
x > 8 .
\]
Setelah $8$ km, perusahaan pertama lebih murah; tepat pada $8$ km keduanya
menarik biaya $16$.
\end{solution}

\begin{solution}{exo:g10:reffunc:9}
Kedua ruas taknegatif, sehingga ketaksamaannya setara dengan ketaksamaan
antara kuadratnya:
\[
\left(\sqrt{a+b}\right)^2 = a + b
\qquad\text{dan}\qquad
\left(\sqrt a + \sqrt b\right)^2 = a + 2\sqrt a \sqrt b + b .
\]
Karena $2\sqrt a\sqrt b \geq 0$, kuadrat kedua paling sedikit sama dengan yang
pertama, dan itu membuktikan $\sqrt{a+b} \leq \sqrt a + \sqrt b$. Kesamaan
berlaku tepat ketika $\sqrt a \sqrt b = 0$, yaitu ketika $a = 0$ atau $b = 0$.
\end{solution}

\begin{solution}{exo:g10:reffunc:10}
\emph{1.} Samakan penyebut ruas kanannya:
\[
2 + \frac{3}{x-1} = \frac{2(x-1) + 3}{x - 1} = \frac{2x + 1}{x - 1}
= f(x).
\]

\emph{2.} Pada $\intoo{1}{+\infty}$, ketika $x$ membesar, $x - 1$ membesar
dan tetap positif, sehingga $\frac{3}{x-1}$ mengecil (\omterm{def:g10:functions:function}{fungsi} kebalikan,
dikalikan konstanta positif $3$), jadi $f(x) = 2 + \frac{3}{x-1}$ mengecil:
$f$ tegas turun pada $\intoo{1}{+\infty}$.
\end{solution}

\begin{solution}{pb:g10:reffunc:1}
\textbf{1.} $d(30) = 3^2 = 9$ m; $d(50) = 25$ m; $d(90) =
81$ m; $d(130) = 169$ m.

\textbf{2.} Melipatduakan $v$ melipatduakan $\frac v{10}$ dan mengalikan
kuadratnya dengan $4$: $d(100) = 100$ m $= 4 \times d(50)$. Laju dua kali
lipat, jarak empat kali lipat --- itulah penskalaan \omterm{def:g10:functions:function}{fungsi} pangkat dua.

\textbf{3.} $\left(\frac{v}{10}\right)^2 = 50$ memberi
$\frac{v}{10} = \sqrt{50}$, jadi $v = 10\sqrt{50} \approx 71$
km/jam. Akar kuadratlah yang menjawab; pembacaannya tidak bermakna ganda
karena laju selalu positif, dan pada
$\intco{0}{+\infty}$ \omterm{def:g10:functions:function}{fungsi} pangkat dua mendaki secara tegas
(\cref{prop:g10:reffunc:sqrt}: satu \omterm{def:g10:functions:function}{prapeta} positif).

\textbf{4.} $d(31) - d(30) = 9.61 - 9 = 0.61$ m, sedangkan
$d(91) - d(90) = 82.81 - 81 = 1.81$ m: tambahan $+1$ km/jam yang sama
menuntut jarak tiga kali lebih banyak pada laju tinggi. Parabolanya tidak
hanya naik, tetapi juga \emph{makin curam} --- laju pendakiannya membesar
seiring $x$.

\textbf{5.} $D(50) = 14 + 25 = 39$ m; $D(130) = 36.4 + 169 =
205.4$ m. Di dalam kota suku afin (reaksi) yang berbagian lebih besar;
di jalan tol suku pangkat dua menghancurkannya --- yang afin
tumbuh mantap, yang kuadrat melesat pergi.

\textbf{6.} $F(0.5) = 240$, $F(1) = 120$, $F(3) = 40$.

\textbf{7.} \omterm{def:g10:functions:function}{Fungsi} kebalikan turun: makin jauh dari
porosnya, makin kecil gaya yang diperlukan --- dengan lengan tuas yang
cukup panjang ($d$ sangat besar), gaya apa pun, sekecil apa pun juga,
menyeimbangkan beban apa pun: kesesumbaran Archimedes hidup di ekor
hiperbolanya, yang mendekati $0$ tanpa pernah mencapainya. Dan $d = 0$
adalah nilai terlarangnya: tanpa lengan, tanpa tuas --- itulah asimtot tegaknya.

\textbf{8.} $\frac{100}{4} = 25$ pada $2$ m,
$\frac{100}{9} \approx 11$ pada $3$ m, $1$ pada $10$ m. Pada jarak
$d$ cahaya itu telah tersebar pada bola berluas $4\pi d^2$
(soal batu nisan Archimedes di jilid sebelumnya): energi yang sama dibagi
luas yang tumbuh seperti $d^2$ --- karena itulah kebalikan \emph{kuadrat}.

\textbf{9.} $60$, $30$, $20$ euro per kepala untuk $n = 10, 20,
30$. Agar $\frac{600}{n} \leq 25$: $n \geq 24$ peserta.
Hiperbolanya menjanjikan bagian yang makin murah ketika $n$ membesar ---
dan menolak untuk pernah mencapai $0$: busnya tidak pernah gratis.

\textbf{10.} $a(x) = \frac{2x + 3}{x} = 2 + \frac3x$:
turun pada $\intoo{0}{+\infty}$ (\omterm{def:g10:functions:function}{fungsi} kebalikan yang tergeser),
mendekati asimtot $y = 2$. Secara ekonomi: menyebarkan biaya
tetap $3$ euro ke lebih banyak poster menekan biaya satuannya turun
menuju $2$ euro bahan yang tak tertekan lagi --- skala ekonomi,
dengan lantai yang keras.

\textbf{11.} Bumi: $a^3 = 1$, $T^2 = 1$. Mars:
$1.52^3 \approx 3.51$ dan $1.88^2 \approx 3.53$. Jupiter:
$5.20^3 \approx 140.6$ dan $11.86^2 \approx 140.7$. Yang diperhatikan
Kepler: $T^2 = a^3$ untuk setiap planet --- satu hukum untuk seluruh
langit.

\textbf{12.} $T = \sqrt{a^3}$: untuk Saturnus
$\sqrt{9.54^3} = \sqrt{868} \approx 29.5$ tahun, terhadap
$29.4$ hasil pengamatan: hukumnya berlaku sampai satu persepuluhan.

\textbf{13.} Asteroid: $T = \sqrt{4^3} = \sqrt{64} = 8$ tahun.
Komet: $a^3 = 27^2 = 729 = 9^3$, jadi $a = 9$ satuan astronomi.

\textbf{14.} Pangkat dua ($T^2$), pangkat tiga ($a^3$) dan akar kuadrat
(untuk menarik $T$). Penskalaannya: mengalikan $a$ dengan $4$ mengalikan
$T$ dengan $4\sqrt4 = 8$ --- seperti yang dibenarkan asteroid tadi ($a$
empat kali punya Bumi, $T$ delapan kali).

\textbf{15.} Sebuah tabel bilangan, yang dicocokkan dengan bentuk
segelintir \omterm{def:g10:functions:function}{fungsi} acuan, menghasilkan sebuah hukum alam semesta
tujuh puluh tahun sebelum seorang pun tahu \emph{mengapa} hukum itu
berlaku --- mengenali \omterm{def:g10:functions:function}{fungsinya} adalah separuh dari ilmu.

\textbf{16.} Pada $\intoo{0}{1}$: $x^2 < x < \sqrt x$; pada
$\intoo{1}{+\infty}$: $\sqrt x < x < x^2$. Pemeriksaannya: di
$x = 0.25$: $0.0625 < 0.25 < 0.5$; di $x = 4$:
$2 < 4 < 16$.

\textbf{17.} Untuk $x \geq 0$, kedua ruas bernilai $\geq 0$, jadi
$\sqrt x > x \iff x > x^2 \iff x(1 - x) > 0 \iff
x \in \intoo{0}{1}$.

\textbf{18.} Di $x = 0.5$: $x^3 = 0.125 < x^2 = 0.25 < x =
0.5$. Di $x = 2$: $2 < 4 < 8$: urutannya terbalik. Perpotongan
$x^2$ dan $x^3$: $x^3 - x^2 = x^2(x - 1) = 0$: di $x = 0$ dan
$x = 1$.

\textbf{19.} $10t^2 = 100\sqrt t$ memberi $t^2 = 10\sqrt t$;
setelah dikuadratkan, $t^4 = 100t$, jadi $t^3 = 100$ dan
$t = \sqrt[3]{100} \approx 4.64$ tahun --- sekitar $4$ tahun dan
$8$ bulan. (Kedua skema lalu membayar sekitar $215$ euro.) Moralnya:
akar berlari kencang di awal, pangkat memenangi setiap maraton.

\textbf{20.} Pengereman: parabola yang \emph{makin curam} adalah
bahaya laju tinggi. Tuas: \emph{asimtot tegak} di $0$
dan ekor yang panjang adalah keuntungan sang montir. Lampu:
$d^2$ dari \emph{bola yang menyebar} menentukan peredupannya. Kepler:
\emph{pangkat} $\frac32$ adalah tempo tata surya.
Tabungan A: akar yang \emph{mendatar} adalah nasib penabung yang sabar.
Lima \omterm{def:g10:functions:graph}{grafik}, lima hukum: sungguh, itulah abjad alam.

\end{solution}
